\documentclass[aspectratio=169,11pt,spanish]{beamer}
\input{../../../_preambulo_beamer.tex}
\input{../../../_comandos_beamer.tex}
\title[L05 · 40 minutos]{Teoría de conjuntos y espacios muestrales}
\begin{document}
\shorthandoff{"}
\begin{frame}{L05 · Un despacho, tres características}
\framesubtitle{00–02 min · Caso}
Un despacho registra \textbf{turno} (Día/Noche), \textbf{servicio} (Estándar/Exprés) y \textbf{estado} (A tiempo/Retrasado).

\medskip
\textbf{Pregunta:} ¿qué despachos pueden ser nocturnos o retrasados?

\medskip
Trabajaremos con ocho tipos posibles y frecuencias de 480 despachos \textbf{sintéticos}. No representan a una empresa real.
\end{frame}
\begin{frame}{Un resultado es una tupla; $S$ reúne las posibilidades}
\framesubtitle{02–07 min · Modelo}
\[
\omega=(\text{turno},\text{servicio},\text{estado}),\qquad S=T\times V\times E.
\]
$T=\{\text{Día},\text{Noche}\}$, $V=\{\text{Estándar},\text{Exprés}\}$ y $E=\{\text{A tiempo},\text{Retrasado}\}$.
\[
|S|=2\cdot2\cdot2=8.
\]
\textbf{Comprueba:} ¿cuáles son las dos tuplas de Noche, Exprés?

\medskip
\textbf{Respuesta:} (Noche, Exprés, A tiempo) y (Noche, Exprés, Retrasado).
\end{frame}
\begin{frame}{Notebook · Enumerar las ocho tuplas}
\framesubtitle{07–10 min · Celda 1}
\texttt{product(TURNOS, SERVICIOS, ESTADOS)} construye $S$ en el orden de las categorías.

\medskip
\textbf{Ejecuta la celda 1.} Comprueba que aparecen ocho tuplas distintas. La primera es (Día, Estándar, A tiempo); la última es (Noche, Exprés, Retrasado).

\medskip
Una tupla identifica un \textbf{tipo de resultado}, no un despacho individual.
\end{frame}
\begin{frame}{Ocho posibilidades; 480 registros}
\framesubtitle{10–14 min · Celda 2}
La celda 2 asigna a las ocho tuplas sus frecuencias sintéticas, generadas por el script extenso con semilla 42.
\[
\sum_{\omega\in S}n(\omega)=480,\qquad |S|=8.
\]
Por ejemplo, (Día, Estándar, A tiempo) aparece 196 veces.

\medskip
\textbf{Comprueba:} si aparece otro despacho de ese tipo, ¿cambia $|S|$?

\medskip
\textbf{Respuesta:} no; aumenta $n(\omega)$, no el número de posibilidades.
\end{frame}
\begin{frame}{Los eventos seleccionan tipos de resultado}
\framesubtitle{14–18 min · Celda 3}
\[
A=\{\omega\in S:\text{turno = Noche}\},\qquad
B=\{\omega\in S:\text{estado = Retrasado}\}.
\]
$A$ y $B$ son subconjuntos de $S$, con cuatro tipos cada uno.

\medskip
\textbf{Antes de ejecutar:} ¿pertenece (Noche, Exprés, Retrasado) a ambos?

\medskip
Sí. La misma tupla puede satisfacer dos condiciones.
\end{frame}
\begin{frame}{Intersección, unión y complemento}
\framesubtitle{18–23 min · Predicción}
\[
\begin{array}{ll}
A\cap B &: \text{Noche y Retrasado}\quad |A\cap B|=2,\\[3pt]
A\cup B &: \text{Noche o Retrasado}\quad |A\cup B|=6,\\[3pt]
A^c=S\setminus A &: \text{Día}\quad |A^c|=4.
\end{array}
\]
La \textbf{o} es inclusiva: una tupla que cumple ambas condiciones entra una sola vez en $A\cup B$.

\medskip
\textbf{Comprueba:} $|A\cup B|=4+4-2=6$.
\end{frame}
\begin{frame}{Notebook · Tipos posibles y registros observados}
\framesubtitle{23–26 min · Celda 3}
\centering
\begin{tabular}{lrr}
\toprule Evento & Tipos posibles & Registros sintéticos\\ \midrule
$A$ & 4 & 172\\
$B$ & 4 & 86\\
$A\cap B$ & 2 & 45\\
$A\cup B$ & 6 & 213\\
$A^c$ & 4 & 308\\ \bottomrule
\end{tabular}
\par\medskip\raggedright
\textbf{Ejecuta la celda 3.} ¿Por qué 45 no es la cardinalidad de $A\cap B$?

\medskip
Porque 45 cuenta despachos observados; la intersección contiene dos tipos.
\end{frame}
\begin{frame}{Disyunción y exhaustividad son preguntas distintas}
\framesubtitle{26–30 min · Definiciones}
Un par es \textbf{disjunto} si $X\cap Y=\varnothing$. Es \textbf{exhaustivo} si $X\cup Y=S$.

\medskip
Para $A$ (Noche) y $B$ (Retrasado):
\[
|A\cap B|=2,\qquad |A\cup B|=6.
\]
No son disjuntos ni exhaustivos.

\medskip
\textbf{Pregunta:} ¿qué dos tipos quedan fuera de $A\cup B$?

\medskip
Los de Día y A tiempo, uno por cada servicio.
\end{frame}
\begin{frame}{Notebook · Un par que sí particiona $S$}
\framesubtitle{30–33 min · Celda 4}
Sea $C$ = Estándar y $D$ = Exprés.
\[
C\cap D=\varnothing,\qquad C\cup D=S.
\]
\textbf{Ejecuta la celda 4.} Contrasta $(A,B)$ con $(C,D)$.

\medskip
\centering
\begin{tabular}{lcc}
\toprule Par & Disjunto & Exhaustivo\\ \midrule
$A,B$ & No & No\\
$C,D$ & Sí & Sí\\ \bottomrule
\end{tabular}
\end{frame}
\begin{frame}{Integra · Una consulta nueva}
\framesubtitle{33–37 min · Celda 5}
Define $F$ = Exprés y Retrasado.

\medskip
\textbf{Antes de ejecutar:} predice $|F|$, los registros sintéticos de $F$, $|F\cap A^c|$ y $|F\cup A^c|$.

\medskip
Pista: $A^c$ es el conjunto de resultados de Día.

\medskip
Ejecuta la celda 5 y explica cada número en su unidad: tipo posible o registro observado.
\end{frame}
\begin{frame}{Comprueba la respuesta}
\framesubtitle{37–39 min · Interpretación}
\[
|F|=2,\qquad n(F)=40,\qquad
|F\cap A^c|=1,\qquad |F\cup A^c|=5.
\]
Una tupla es Exprés, Retrasado y Día. La otra es Exprés, Retrasado y Noche.

\medskip
\textbf{Salida de clase:} explica por qué $|F|$ y $n(F)$ responden preguntas distintas.
\end{frame}
\begin{frame}{Para seguir estudiando}
\framesubtitle{39–40 min · Cierre}
El notebook extenso desarrolla la simulación de los 480 despachos, las figuras y más ejercicios. Los tres scripts muestran cómo se generan y analizan los registros.

\medskip
\textbf{Puente a L06:} cuando el espacio muestral crece, enumerarlo a mano deja de ser práctico. Necesitaremos reglas de conteo.
\end{frame}
\end{document}
